\documentclass[aspectratio=169,10pt]{beamer}
\usetheme[numbering=fraction,progressbar=frametitle]{metropolis}
\usepackage{fontspec}
\setsansfont{Latin Modern Sans}
\usepackage{amsmath,amssymb,mathtools,graphicx,booktabs}
\usepackage{polyglossia}
\setdefaultlanguage{english}
\setbeamertemplate{navigation symbols}{}
\setbeamerfont{frametitle}{size=\large}
\setbeamerfont{normal text}{size=\normalsize}
\setlength{\parskip}{.45em}
\hypersetup{pdfauthor={},pdfsubject={Lecture slides},pdfkeywords={},pdftitle={Real numbers and intervals}}
\providecommand{\N}{\mathbb{N}}
\providecommand{\Z}{\mathbb{Z}}
\providecommand{\D}{\mathbb{D}}
\providecommand{\Q}{\mathbb{Q}}
\providecommand{\R}{\mathbb{R}}
\providecommand{\C}{\mathbb{C}}
\providecommand{\vect}[1]{\overrightarrow{#1}}
\providecommand{\vi}{\vec{\imath}}
\providecommand{\vj}{\vec{\jmath}}
\providecommand{\vk}{\vec{k}}
\providecommand{\norm}[1]{\left\lVert #1\right\rVert}
\providecommand{\vc}[2]{\begin{pmatrix}#1\\#2\end{pmatrix}}
\providecommand{\deter}[2]{\det\left(#1,#2\right)}
\providecommand{\abs}[1]{\left\lvert #1\right\rvert}
\providecommand{\dd}{\,\mathrm{d}}
\providecommand{\e}{\mathrm{e}}
\providecommand{\Card}{\operatorname{Card}}
\providecommand{\E}{\operatorname{E}}
\providecommand{\V}{\operatorname{V}}
\providecommand{\degre}{^{\circ}}
\providecommand{\un}[1]{\,\mathrm{#1}}
\providecommand{\sci}[2]{#1\times10^{#2}}
\providecommand{\pgcd}{\operatorname{PGCD}}
\providecommand{\ppcm}{\operatorname{PPCM}}
\title{Real numbers and intervals}
\subtitle{Lecture and worked examples}
\date{}
\author{}
\begin{document}
\begin{frame}[plain]
\vfill
{\usebeamerfont{title}\inserttitle\par}
\medskip
{\usebeamerfont{subtitle}\insertsubtitle\par}
\medskip
{\usebeamercolor[fg]{progress bar}\rule{\textwidth}{.5pt}}
\vfill
\end{frame}
\begin{frame}{Sets of numbers}
\begin{columns}[T,onlytextwidth]
\begin{column}{.57\textwidth}
\begin{itemize}

\item \(\N\): natural numbers; \(\Z\): integers.

\item \(\D\): decimal numbers, of the form \(a/10^n\) with \(a\in\Z\), \(n\in\N\).

\item \(\Q\): rational numbers, of the form \(a/b\) with \(a,b\in\Z\), \(b\ne0\).

\item \(\R\): real numbers, represented by the abscissae of the points of a graduated line. A real number that is not rational is \textbf{irrational}.

\end{itemize}

\[\N\subsetneq\Z\subsetneq\D\subsetneq\Q\subsetneq\R.\]
\end{column}
\begin{column}{.4\textwidth}
\centering
\includegraphics[width=\linewidth,height=4.5cm,keepaspectratio]{figures/seconde-04-1-en.pdf}
\end{column}
\end{columns}
\end{frame}
\begin{frame}{Sets of numbers}
Give the smallest of the five sets containing each of the numbers: \(-95\); \(7.29\); \(5/3\); \(\sqrt{25}\); \(\sqrt2\); \(\pi\); \(0.0002\).
\end{frame}
\begin{frame}{Worked answer: Sets of numbers}
\(-95\in\Z\); \(7.29\in\D\); \(5/3\in\Q\); \(\sqrt{25}=5\in\N\); \(\sqrt2,\pi\in\R\setminus\Q\); \(0.0002=2/10^4\in\D\). The irrationality of \(\pi\) is accepted here.
\end{frame}
\begin{frame}{Two proofs by contradiction}
To prove a statement, assume its negation, deduce a contradiction, then conclude that the statement is true.
\end{frame}
\begin{frame}{Two proofs by contradiction}
\(\dfrac13\) is not a decimal number.
\end{frame}
\begin{frame}{Two proofs by contradiction}
Suppose \(1/3=a/10^n\) with \(a\in\Z\), \(n\in\N\). Then \(10^n=3a\) would be a multiple of \(3\). Yet \(10\equiv1\) modulo \(3\), so \(10^n\) leaves remainder \(1\) on division by \(3\) (the divisibility rule for \(3\) can also be used). Contradiction. Hence \(1/3\notin\D\).
\end{frame}
\begin{frame}{Two proofs by contradiction}
\(\sqrt2\) is irrational.
\end{frame}
\begin{frame}{Two proofs by contradiction}
Suppose \(\sqrt2=p/q\) with \(p,q\in\N\), \(q\ne0\), and \(p/q\) in lowest terms. Then \(p^2=2q^2\), so \(p^2\) is even. If \(p\) were odd, its square would be odd: so \(p\) is even. Write \(p=2k\). We get \(4k^2=2q^2\), then \(q^2=2k^2\); so \(q\) is even in turn. Thus \(2\) divides \(p\) and \(q\), contradicting irreducibility. Hence \(\sqrt2\notin\Q\).
\end{frame}
\begin{frame}{Intervals of the real line}
For two real numbers \(a<b\), intervals describe the real numbers lying between two endpoints, with or without the endpoints. A square bracket includes the endpoint, a parenthesis excludes it. The symbols \(-\infty,+\infty\) are not real numbers and always take a parenthesis.
\end{frame}
\begin{frame}{Intervals of the real line}
The strictly negative real numbers form \((-\infty,0)\): \(0\) is excluded. \([a,a]=\{a\}\); \((a,a)=\varnothing\).
\end{frame}
\begin{frame}{Intervals of the real line}
Write as an interval then represent: \(2\leq x\leq5\); \(x>4\); \(-6<x\leq2\); \(x<0\).
\end{frame}
\begin{frame}{Worked answer: Intervals of the real line}
\([2,5]\); \((4,+\infty)\); \((-6,2]\); \((-\infty,0)\). On the number line, mark the endpoints, then show each one as included or excluded.
\end{frame}
\begin{frame}{Intersection and union of intervals}
\(\,I\cap J\) corresponds to the condition “belonging to \(I\) \textbf{and} to \(J\)”; \(I\cup J\) corresponds to “belonging to \(I\) \textbf{or} to \(J\)”. A union of intervals is not always an interval.
\end{frame}
\begin{frame}{Intersection and union of intervals}
With \(I=[0,12]\) and \(J=[3,20]\), determine \(I\cap J\) and \(I\cup J\).
\end{frame}
\begin{frame}{Worked answer: Intersection and union of intervals}
\(I\cap J=[3,12]\) and \(I\cup J=[0,20]\).
\end{frame}
\begin{frame}{Intersection and union of intervals}
Complete, justifying with representations on a number line.
\end{frame}
\begin{frame}{Worked answer: Intersection and union of intervals}
For each intersection, keep the points common to both representations; for each union, keep all the points of at least one. In the last row, \(0\) belongs to neither interval: the union cannot be replaced by \((-2,3]\).
\end{frame}
\begin{frame}{Absolute value and distance}
\(|x|\) is the distance from the real number \(x\) to \(0\); \(|x|\geq0\) and
\[|x|=\begin{cases}x&\text{if }x\geq0,\\-x&\text{if }x<0.\end{cases}\]
The distance between \(a\) and \(b\) is \(|b-a|\). For \(a\in\R\) and \(r\geq0\):
\[|x-a|\leq r\iff x\in[a-r,a+r],\qquad |x-a|<r\iff x\in(a-r,a+r).\]
\end{frame}
\begin{frame}{Absolute value and distance}
Calculate \(|3.2|\), \(|-1.6|\), \(|\sqrt3-2|\), then the distances between \(2\) and \(5\), \(2\) and \(-5\), \(-2\) and \(-5\).
\end{frame}
\begin{frame}{Worked answer: Absolute value and distance}
\(3.2\); \(1.6\); \(2-\sqrt3\) since \(\sqrt3<2\). The distances are \(3\), \(7\) and \(3\).
\end{frame}
\begin{frame}{Absolute value and distance}
Solve: \(|x-7/3|=14\); \(|x-2|\leq3\); \(|x+3|<5\).
\end{frame}
\begin{frame}{Worked answer: Absolute value and distance}
For the first equation, \(x-7/3=14\) or \(x-7/3=-14\), so \(S=\{-35/3,49/3\}\). Then \(|x-2|\leq3\iff x\in[-1,5]\). Finally \(|x+3|=|x-(-3)|<5\iff x\in(-8,2)\).
\end{frame}
\begin{frame}{Absolute value and distance}
Write \(x\in[4,10]\) as an inequality with an absolute value.
\end{frame}
\begin{frame}{Worked answer: Absolute value and distance}
The centre is \((4+10)/2=7\) and the radius \((10-4)/2=3\): \(|x-7|\leq3\).
\end{frame}
\begin{frame}{Bounds and approximations}
To bound \(x\) is to give \(a\leq x\leq b\); the width is \(b-a\). An approximation from below is less than or equal to \(x\); an approximation from above is greater than or equal to it. Rounding chooses the nearest value at the required precision.
Significant figures are counted from the first nonzero digit; the displayed precision must suit the context.
\end{frame}
\begin{frame}{Bounds and approximations}
Bound \(\sqrt2\) between two consecutive hundredths. Give its rounding to the nearest hundredth. Round \(0.004567\) to three significant figures, then \(12345\) to three significant figures in scientific notation.
\end{frame}
\begin{frame}{Worked answer: Bounds and approximations}
\(1.41^2=1.9881<2<2.0164=1.42^2\), so \(1.41<\sqrt2<1.42\) (width \(0.01\)). The rounding to the nearest hundredth is \(1.41\). The other roundings are \(0.00457\) and \(1.23\times10^4\).
\end{frame}
\begin{frame}[fragile]{Python}
\small
\begin{verbatim}
def encadre_racine2(p):
    # p entier naturel ; calculs entiers exacts
    q = 10**p
    k = 0
    while (k+1)**2 <= 2*q*q:
        k = k+1
    return k/q, (k+1)/q
\end{verbatim}
\end{frame}
\begin{frame}{Bounds and approximations}
Explain why this function bounds \(\sqrt2\). What does it return for \(p=2\)? Why does it terminate?
\end{frame}
\begin{frame}{Worked answer: Bounds and approximations}
The loop looks for the largest integer \(k\) such that \(k^2\leq2q^2\). On exit, \(k^2\leq2q^2<(k+1)^2\), so \(k/q\leq\sqrt2<(k+1)/q\). For \(p=2\), it returns \((1.41,1.42)\). Since \((2q)^2>2q^2\), the counter eventually reaches a value that makes the test false. The final divisions are floating-point approximations; the square tests are exact.
\end{frame}
\end{document}
